\begin{exercicebox}[\hfill \it{6 points}]
    Les surfaces de 30 terrains d'un lotissement se répartissent selon le
    tableau suivant.

    \begin{center}
    \begin{tabular}{|c|c|c|c|c|c|c|c|}
        \hline
        \bb{Surface des terrains en (m2)} & 265 & 330 & 430 & 550 & Total  \\ \hline
        \bb{Nombre de terrains}  & 9 &  & 6 & 3 &   \\ \hline
        &&&&&\\
        \bb{Fréquence sous forme fractionnaire} & \hspace{1.3cm} & \hspace{1.3cm} &
        \hspace{1.3cm} &\hspace{1.3cm}  & \\ 
        &&&&&\\\hline
        &&&&&\\
        \bb{Fréquence en \%}  &  &  &  &  &    \\ 
        &&&&&\\
        \hline
    \end{tabular}
    \end{center}
   
    \begin{enumerate}
        \item \bb{Compléter} le tableau directement sur le polycopié. (On ne demande
            pas de justifier)
        \item Alix affirme "Moins du quart des terrains ont une surface
            supérieure à 400m2.". Est-ce exact ? \bb{Justifier.}
    \end{enumerate}
\end{exercicebox}
